\chapter{Metrics Analysis}
\label{ch:metricsAnalysis}
\setcounter{section}{-1} % Imposta il contatore a -1 così la prossima sezione sarà 0

\begin{comment}
In this chapter, we evaluate the effectiveness of the metrics in \cref{ch:metrics}
in detecting the malicious versions of the packages of the four attacks considered in \cref{sec:malwareDataset}.

To validate each metric, we follow a three-step analysis:
\begin{enumerate}
    \item An analysis of the metric's behavior considering the real-world NPM package dataset \cref{sec:goodwareDataset},
    \item An investigation of the anomalous values found therein,
    \item A comparison between the real-world package dataset and the malicious package dataset \cref{sec:malwareDataset}.
\end{enumerate}


As mentioned in \cref{ch:metrics}, these metrics were defined based on the main characteristics
of one or more of the four attacks.
Therefore, we do not expect every metric to detect all four attacks.
Rather, we expect each metric to detect the specific attack or attacks for which it was defined.
\end{comment}

\section{Preface to the analysis}
\label{sec:prefaceAnalysis}

\subsection{Goodware Dataset Scenario}

Per ogni metrica, sono stati presi in analisi i pacchetti contenuti nel dataset goodware, definito in \cref{sec:goodwareDataset}. 
In particolare, tra i 1000 pacchetti più scaricati da NPM, sono stati considerati solo quelli che hanno almeno 20 versioni disponibili,
ovvero 629 pacchetti.

Come descritto nei paragrafi di analisi, tutti i pacchetti non significativi per tale metrica, 
verranno esclusi dai successivi grafici di analisi.

\subsection{Outlier Analysis}

Escludendo la prima metrica, trattata nel capitolo \cref{sec:analysisFileTypes},
tutti i grafici delle metriche si leggono in questo modo.

I grafici di analisi illustrano 
l'andamento nelle 20 versioni
della metrica presa in analisi.


L'asse delle ascisse riporta le versioni dei pacchetti dalla meno recente (-19) all'ultima disponibile (0), mentre 
l'asse delle ordinate rappresenta il numero delle occorrenze della metrica.

Per meglio rappresentare la popolazione eterogenea dei valori delle ordinate, viene adottata la scala logaritmica.


Le altezze delle barre nere verticali rappresentano gli intervalli di confidenza, 
ovvero gli intervalli in cui risiedono il 95\% della popolazione dei pacchetti per una determinata versione.
I pacchetti che si posizionano all'interno di uno di questi intervalli
hanno un numero di valori compreso 
tra il valore minimo (0) e il 95° percentile della popolazione per quella versione.

I pallini grigi invece rappresentano i valori outlier,
ovvero il restante 5\% della popolazione che si posiziona al di fuori dell'intervallo di confidenza.
Un pacchetto assume valori outlier per una versione quando il numero di valori in quella versione supera il valore del 95° percentile
di quell'intervallo di confidenza.


Per ogni metrica il grafico sopra descritto vengono riportati 2 grafici distinti.
Il primo illusta l'andamento nelle 20 versioni
della metrica presa in analisi del dataset goodware, definito nel capitolo x;
il secondo mostra l'andamento nelle 20 versioni
della metrica presa in analisi confrontando il dataset goodware con il dataset Malware, 
definito nel capitolo x.

Eccezzioni: nell'ambito di ogni metrica sono stati eventualmente evidenziati pacchetti con valori 
fuori scala, ossia un numero di occorrenze tale da 
essere per più o tutte le versioni un outlier o il limite superiore di un intervallo di confidenza.


Questi pacchetti, e i loro relativi valori outlier, hanno 
un comportamento atipico rispetto alla maggior parte della popolazione,
motivo per cui sono stati evidenziati nel primo grafico di ciascuna metrica con linee colorate
e successivamente analizzati singolarmente
per comprenderne la natura e le motivazioni che hanno portato a questi valori esadecimali unici così elevati.



Nel primo grafico di ogni metrica è stata indicata con linea tratteggiata in rosso la mediana
calcolata sui pacchetti del dataset Goodware.

A differenza della media, la mediana non è influenzata dai valori outlier e  
non distorce quindi il reale comportamento della metrica.

E' stata inserita la mediana per evidenziare questa forte assimmetria tra i pacchetti,
in quanto esistono pacchetti con un numero di valori bassi, e
pacchetti con un numero di valori che superano di diversi ordini di grandezza
quelli della maggior parte della popolazione.
Tali comportamenti veranno analizzati singolarmente per ciascuna metrica.


\subsection{Comparison Malware dataset vs. Goodware dataset and Results for each 4 attacks}

Una volta compresa la natura di questi valori outlier, 
viene riportato il secondo grafico per ciascuna metrica,
per verificare se le versioni malevole dei pacchetti, relative ai 4 specifici attacchi considerati
nel capitolo x, siano individuabili con tale metrica.

Le versioni malevole, affinchè siano individuabili, devono 
posizionarsi tra i valori outlier, ovvero superare 
il valore del 95° percentile dell'intervallo di confidenza per quella versione.

Se invece, le versioni malevole assumono un numero di 
occorrenze tale da posizionarsi dentro agli
intervalli di confidenza, ovvero tra il valore minimo e il valore del 95° percentile;
la metrica non sarà indicativa per rilevarle,
in quanto risulterebbero indistinguibili da quelle legittime.


Il secondo grafico mantiene gli stessi intervalli di confidenza 
e outlier del primo grafico, 
calcolati sempre sugli stessi pacchetti del dataset goodware, nelle 20 versioni.

Non sono più presenti né le linee colorate che rappresentano i pacchetti evidenziati precedentemente,
né la mediana, in quanto non più necessari per lo scopo di questo grafico.

Nel secondo grafico, le linee colorate con diverse tipologie di tratteggio
rappresentano i pacchetti del dataset Malware, ossia i pacchetti 
compromessi da uno dei 4 specifici attacchi,
descritti nel capitolo x.
Le croci rosse del secondo grafico rappresentano le versioni malevole.

\newpage

\section{Analysis Metric 1: File Types}
\label{sec:analysisFileTypes}

\subsection{Goodware Dataset Scenario}
\label{sec:goodwareFileTypes}


\subsection{Outlier Analysis}
\label{sec:outlierFileTypes}



\subsection{Comparison Malware dataset vs. Goodware dataset and Results for each 4 attacks}




\newpage

\begin{comment}
% ------------------------- QUESTO E' DA DECOMMENTARE
Rarely, some packages may not have any JavaScript or TypeScript code files at all, 
and this depends on the purpose of the package.
For example, there are packages whose resource and utility lie solely in the data inside \texttt{.json} files.
Notable examples include \texttt{node-\allowbreak releases}, which offers a list of Node.js release versions, 
and \texttt{spdx-\allowbreak license-\allowbreak ids}, which provides the official list of \ac{SPDX} licenses. 
Developers use these packages as dependencies to get the current data without hard-coding it manually.

Alternatively, some packages may have no code at all because they are deprecated and maintained solely as placeholders, 
like recently \texttt{@types/\allowbreak uuid} package.


As can be observed better in \nameref{ch:analysis}, other files can be found in tarballs.

The \texttt{html} files, after careful analysis,
represent mainly test files, code coverage reports, and demos, such in the \texttt{jwt-decode} package;
or presentation pages of the package, documentation, and licenses, such in the \texttt{sax} and \texttt{pino} packages.

Instead, the \texttt{css} files are either used as test fixtures (e.g., example data for testing), 
as in the case of the \texttt{convert-source-map} package;
or simply as styles for HTML-based tests and demos, as in the \texttt{loglevel} package.

% ripeto, NON TROVO UNA FONTE DECENTE CHE DICE QUESTO
However, in modern packaging practices,
non essential assets are often removed to keep the final product as lightweight as possible.

Moreover, like build artifacts, these files are not always located in specific directories,
and it is not possible to determine with certainty if a file is test file or not.
\end{comment}


\begin{comment}
Analyzing the dataset also uncovers the historical development of structure in tarballs.
In the past, tarballs were more likely to include configuration files from the GitHub repository, 
for instance, .gitattributes, .gitmodules, and .gitignore (see the json5 package for an example). 
Developers gradually shifted towards omitting such files to make the packages lighter and only essential files would get to the end users.

The same decline pattern is observed in the use of .html files which fell from an initial 3\% (17 packages) to around 1\% (7 packages) 
in the versions of the packages released later. 
A review of these files reveals that they mostly take the form of test suites, code coverage, or demos (e.g., in the jwt, decode package), 
as well as documentation or license pages (e.g., in sax and pino).

Lastly, .css files are typically used either as test fixtures (i.e., example data for testing), 
as in the case of the convert, source, map package, or as styling for HTML, based tests and demos, as in loglevel.
Keeping with the trend of modern packaging, the final product is made as relaxed as possible, non, essential assets are removed.
\end{comment}



% ------------------------------------------
\section{Analysis Metric 2: Package Size}

\subsection{Goodware Dataset Scenario}


\subsection{Outlier Analysis}



\subsection{Comparison Malware dataset vs. Goodware dataset and Results for each 4 attacks}

\newpage



% ------------------------------------------
\section{Analysis Metric 3: Unique Hexadecimal values (no duplicate values)}

\subsection{Goodware Dataset Scenario}

In questo capitolo si analizza
l'andamento nelle 20 versioni
della metrica relativa al numero di valori esadecimali unici, 
ovvero non considerando valori duplicati,
presenti nei file plain text dei pacchetti, come definito nel capitolo metrica x.


Nel grafico \ref{fig:hex_unique_plot1} sono stati esclusi dai 629 pacchetti del dataset Goodware
i 522 pacchetti che non hanno avuto alcun valore esadecimale positivo in nessuna delle 20 versioni
e pertanto nessun valore anomalo da analizzare.

I 107 pacchetti restanti, considerati nel grafico \ref{fig:hex_unique_plot1}, sono quindi 
quelli che hanno avuto almeno un valore esadecimale positivo in almeno una delle 20 versioni.

%Ti segnalo anche che nella tesi ti dovrai riferire ai grafici come “figure”. E già che ci siamo: “a figure” 
%generica in minuscolo, mentre “Figura x.y” in maiuscolo perché è proprio una precisa. 

\subsection{Outlier Analysis}


\begingroup % Inizia un gruppo per confinare il cambiamento
\renewcommand{\thefigure}{7.3.b} % Sostituisci x.y con il numero che desideri
\begin{figure}[H]
    \centering
    \includegraphics[width=\textwidth]{imgs/hex_unique_plot1.pdf}
    \caption{Real-world packages analysis for the unique hexadecimal values (no duplicates) metric}
    \label{fig:hex_unique_plot1}
\end{figure}
%``unique hexadecimal values (no duplicates)'' metric
\endgroup % Chiude il gruppo: la prossima figura tornerà alla numerazione automatica


Con riferimento alle definizioni di cui al capitolo \ref{sec:prefazioneAnalisi}:

Nell'ambito dei 107 pacchetti, 
sono stati evidenziati 10 pacchetti che presentano caratteristiche interessanti, 
evidenziate con linee colorate.
I restanti pacchetti sono dentro agli intervalli di confidenza, e pertanto poco significati per 
l'analisi di questo scenario.


Questi 10 pacchetti hanno un numero di valori esadecimali unici tale da 
essere per più o tutte le versioni un outlier o il limite superiore di un intervallo di confidenza,
la loro singola analisi viene riportata nel seguito.


Nel grafico \ref{fig:hex_unique_plot1} la mediana è
calcolata sui 107 pacchetti analizzati e rappresentata dalla linea tratteggiata in rosso. 
In ogni versione rimane costante a 1, indicando che almeno il 50\% della popolazione analizzata ha al massimo un unico valore esadecimale.


E' stata inserita la mediana per evidenziare questa forte assimmetria tra i pacchetti,
in quanto esistono pacchetti con un numero di valori esadecimali unici basso, e
pacchetti con un numero di valori esadecimali unici che superano di diversi ordini di grandezza
quelli della maggior parte della popolazione; per vari motivi che discuteremo dopo.


L'obiettivo principale quindi è studiare la natura dei pacchetti che assumono valori esadecimali outlier.

Dal grafico \ref{fig:hex_unique_plot1} si nota che i pacchetti evidenziati
hanno 3 tipologie di trend: crescente, decrescente o stabile.

\begin{itemize}
    \item \textbf{Pacchetti con trend crescente:}
I pacchetti che hanno avuto un incremento considerevole da una versione alla successiva sono: crypto-js, xml2js e polished.

    \begin{itemize}
        \item \textbf{crypto-js:} 
Questo pacchetto è legato alla crittografia, ovvero è una libreria JavaScript 
che implementa primitive crittografiche quali algoritmi di cifratura, %(es. AES(Advanced Encryption Standard))
firme digitali ecc. %hash (es. SHA\-256),

Rappresentato nel grafico \ref{fig:hex_unique_plot1} dalla linea continua nella tonalità del rosa-rosso;
il pacchetto nelle prime 3 versioni di 13 anni fa dalla versione 3.1.2 alla versione 3.1.2-2
(versione -19 alla -17 nel grafico \ref{fig:hex_unique_plot1}), 
usava solo valori decimali per implementare le proprie primitive crittografiche.

Non presentando quindi nessun valore esadecimale, sul grafico \ref{fig:hex_unique_plot1} 
si posizionava a 0.

In un aggiormaneto successivo, il 3.1.2-3 (-16 nel grafico \ref{fig:hex_unique_plot1}), 
gli sviluppatori hanno convertito i valori 
in formato esadecimale per garantire una rappresentazione dei dati più compatta. %dati binari e più fedele.
%riducendo in questo modo il rischio di errori di formattazione.

L'aggiornamento ha incrementato il numero di valori esadecimali unici arrivando a 546,
rappresentando un valore outlier, anche per le versioni successive.

L'incremento invece visibile nel grafico \ref{fig:hex_unique_plot1} dalla versione -1 alla versione 0, 
è dato dall'introduzione di un nuovo algoritmo di crittografia, Blowfish, nella versione 4.2.0.

In questo caso il numero di valori esadecimali unici è arrivato a 1588, 
rappresentando l'outlier più estremo per l'ultima versione (0 nel grafico \ref{fig:hex_unique_plot1}).

        \item \textbf{xml2js:}

Rappresentato nel grafico \ref{fig:hex_unique_plot1} dalla linea continua nella leggera tonalità di rosso;
il pacchetto xml2js ha avuto un incremento considerevole di valori 
esadecimali unici nella versione 0.6.1 (versione -1 nel grafico \ref{fig:hex_unique_plot1}),
passando da 0 a 96 valori esadecimali unici.

Questo incremento è dovuto al fatto che a partire da questa versione, gli sviluppatori 
hanno incluso nel pacchetto uno script JavaScript generato.
Il file era originariamente scritto in OCaml, 
ma per garantirne l'esecuzione in ambienti Node.js, 
il codice è stato convertito automaticamente in JavaScript tramite l'uso
del compilatore js\_of\_ocaml.

Per motivi di efficienza, il processo di compilazione produce uno script che fa largo uso di costanti esadecimali.
%, come ad esempio per gestire in modo più efficiente la rappresentazione dei dati binari e del runtime originale.

Il file JavaScript generato è rimasto incluso nel tarball del pacchetto
anche per la versione successiva 0.6.2 (versione 0 nel grafico \ref{fig:hex_unique_plot1}).

        \item \textbf{Polished:}
        
Rappresentato nel grafico \ref{fig:hex_unique_plot1} dalla linea continua nella tonalità scura di blu-magenta;
il pacchetto polished è la libreria di utilità per lo styling in JavaScript.
Ha avuto un incremento massiccio di valori esadecimali unici nella versione 4.0.0-beta.2 (versione -15 nel grafico \ref{fig:hex_unique_plot1}),
passando da 0 a 2208 valori esadecimali unici, rappresentando l'outlier più estremo per 4 versioni consecutive.

In questa versione, gli sviluppatori hanno erroneamente inserito 
nel tarball finale distribuito al pubblico, una directory cache .yarn.
La directory conteneva numerosi file, tra cui un file JavaScript bundle generato, 
contenente numerosi valori esadecimali.

Nella versione 4.0.4 (-11 nel grafico \ref{fig:hex_unique_plot1}) è stata rimossa
la cartella dal tarball, come riportato nelle note di rilascio
su GitHub (\url{https://github.com/styled-components/polished/releases?page=2}).
    \end{itemize}

    \item \textbf{Pacchetti con trend decrescente:}
I pacchetti che hanno avuto un decremento considerevole da una versione alla successiva sono: buffer-crc32 e object-hash.

    \begin{itemize}
        \item \textbf{buffer-crc32:}

Rappresentato nel grafico \ref{fig:hex_unique_plot1} dalla linea continua nella tonalità di ciano chiaro;
il pacchetto buffer-crc32 è una libreria che implementa l'algoritmo CRC32 (Cyclic Redundancy Check),
usato principalmente per rilevare errori accidentali nella trasmissione o archiviazione nei dati.

Per rendere il controllo dell'integrità dei dati veloce ed efficiente, l'algoritmo 
fa uso della tabella CRC: una matrice di 256 valori pre-calcolati che evita al processore di 
eseguire calcoli matematici complessi su ogni singolo byte.

Fino alla versione 0.2.13 (-9 nel grafico \ref{fig:hex_unique_plot1}), 
il pacchetto aveva 256 valori esadecimali unici, rappresentanti i valori della tabella CRC.

Dalla versione 1.0.0-RC1 (-8 nel grafico \ref{fig:hex_unique_plot1}), gli sviluppatori hanno deciso 
di ottimizzare lo spazio del pacchetto,
distribuendo nel tarball il file contenente la tabella crc generato, usando 
strumenti di build automatici quali bundler e minificatori.
In questo processo automatico, 
gli strumenti di build scelgono spesso il formato decimale perché più compatto;
motivo per cui sono stati convertiti i valori della tabella da esadecimali in decimali,
lasciando però invariato il valore matematico.

        \item \textbf{object-hash:}

Rappresentato nel grafico \ref{fig:hex_unique_plot1} dalla linea continua nella tonalità di magenta;
il pacchetto object-hash è una libreria che 
genera un hash a partire da oggetti e valori JavaScript.

Dalla versione 1.1.0 fino alla 1.1.2 (-19 a -17 nel grafico \ref{fig:hex_unique_plot1}),
il pacchetto aveva un numero di valori esadecimali pari a 78, che rappresentava 
il limite superiore degli intervalli di confidenza per quelle versioni.

Fino alla versione 1.1.6 (-13 nel grafico \ref{fig:hex_unique_plot1}), 
dentro al tarball era presente un file, object\_hash\_test.js, 
contenente test utili a verificare il corretto 
funzionamento della libreria; facendo uso di valori esadecimali.

% NON TROVO UNA FONTE DECENTE CHE DICE QUESTO 
% neanche questo va bene https://snyk.io/blog/best-practices-create-modern-npm-package/
% forse c'è github acction ma non è questo il caso, https://docs.github.com/en/actions/tutorials/build-and-test-code/nodejs
Secondo le linee guida di [OpenSSF] (\url{https://github.com/ossf/package-manager-best-practices/blob/main/published/npm.md}),
è una best practice moderna includere nel tarball solo i file essenziali per il funzionamento della libreria,
escludendo ad esempio file di test.
Ciò è necessario per ridurre la dimensione del pacchetto e 
migliorare le performace di installazione.

%@misc{openssf_npm_best_practices,
%  author        = {{Open Source Security Foundation (OpenSSF)}},
%  title         = {npm Best Practices Guide},
%  howpublished = {\url{https://github.com/ossf/package-manager-best-practices/blob/main/published/npm.md}},
%  year         = {2023}, % L'anno dell'ultima revisione significativa o archiviazione
%  note         = {Accessed: 2024-05-22}
%}
% oppure
%[Brian] (\url{https://snyk.io/blog/best-practices-create-modern-npm-package/})

Il pacchetto object-hash ha seguito questa best practice, 
e a partire dalla versione 1.1.7 (-12 nel grafico \ref{fig:hex_unique_plot1})
ha ridotto prima il numero di test, e di conseguenza diminuito il numero di valori esadecimali unici da 78 a 5.
Successivamente, dalla versione 2.0.3 (-4 nel grafico \ref{fig:hex_unique_plot1}),
ha rimosso completamente il file di test dal tarball, portando il numero di valori esadecimali unici a 0.

    \end{itemize}

    \item \textbf{Pacchetti con andamento stabile:}
I rimanenti 5 pacchetti hanno avuto valori outlier, o valori che rappresentano il limite superiore degli intervalli di confidenza,
sempre costanti per tutte le 20 versioni, senza incrementi o decrementi significativi.
I pacchetti con questo comportamento sono: node-forge, tweetnacl, chardet, next e aws-sdk.

    \begin{itemize}
        \item \textbf{node-forge e tweetnacl:}
I primi due pacchetti, node-forge e tweetnacl sono altre librerie legate alla crittografia,
e nel grafico \ref{fig:hex_unique_plot1} sono rappresentati rispettivamente dalla linea continua nella tonalità di verde 
e dalla linea continua nella tonalità di blu.
Il primo, node-forge, assume un numero di valori esadecimali unici pari a 406, 
rappresentando un outlier in tutte le versioni.

Il secondo, tweetnacl, assume un numero di valori esadecimali unici pari a 177,
rappresentando il limite superiore degli intervalli di confidenza per le versioni -16 e dalla -11 alla -9
nel grafico \ref{fig:hex_unique_plot1}, e rappresentando un outlier per le versioni dalla -19 alla -17, e dalla -8 alla 0.

        \item \textbf{chardet:}
Il pacchetto chardet è una libreria che rileva la codifica dei caratteri, e per funzionare 
fa uso di tabelle e mappe di byte sotto forma di costanti esadecimali, 
principalmente presenti nel file /lib/encoding/sbcs.js.

Rappresentato nel grafico \ref{fig:hex_unique_plot1} dalla linea continua nella tonalità di arancione,
il pacchetto assume un numero di valori esadecimali unici pari a 972, rappresentando l'outlier più estremo 
per la maggior parte delle versioni; superato solo dal pacchetto polished dalla versione -15 alla versione -12, e
nell'ultima versione 0 dal pacchetto crypto-js.


        \item \textbf{next e aws-sdk:}
Rappresentati nel grafico \ref{fig:hex_unique_plot1} rispettivamente dalla linea continua nella tonalità di magenta 
e dalla linea continua nella tonalità di giallo-verde chiaro;
i pacchetti next e aws-sdk, contengono nei loro rispettivi tarball 
numerosi file JavaScript generati di build, che fanno uso di valori esadecimali.
Questo comportamento è il medesimo 
già visto precedentemente per i pacchetti polished e buffer-crc32.

Nel dettaglio: il pacchetto next, framework full-stack di React,
rappresenta nel grafico \ref{fig:hex_unique_plot1}
il limite superiore degli intervalli di confidenza per le versioni dalla -8 alla -2, 
con un numero di valori esadecimali unici costante a 40;

il pacchetto aws-sdk, libreria ufficiale per interagire con i servizi cloud AWS,
rappresenta nel grafico \ref{fig:hex_unique_plot1} 
un outlier in tutte le versioni, con un numero di valori esadecimali unici di 600.

    \end{itemize}
\end{itemize}



\subsection{Comparison Malware dataset vs. Goodware dataset and Results for each 4 attacks}



\begingroup % Inizia un gruppo per confinare il cambiamento
\renewcommand{\thefigure}{7.3.c} % Sostituisci x.y con il numero che desideri
\begin{figure}[H]
    \centering
    \includegraphics[width=\textwidth]{imgs/hex_unique_plot2.pdf}
    \caption{Packages analysis for the unique hexadecimal values (no duplicates) metric with malicious packages versions}
    \label{fig:hex_unique_plot2}
\end{figure}
\endgroup % Chiude il gruppo: la prossima figura tornerà alla numerazione automatica



Come discusso nel capitolo metriche, ho scelto di considerare il numero di valori esadecimali unici come metrica 
perchè esistono attacchi che fanno uso di tecniche di offuscamento,
e secondo [Moog] la tecnica più comunemente usata 
è quella della sostituzione di variabili e nomi di funzioni con identificatori esadecimali randomici.

Analizziamo quando la metrica è efficace per rilevare questi attacchi, e quando invece non è indicativa.
Il funzionamento dei 4 attacchi presi in analisi sono spiegati in dettaglio nel capitolo threat modeling.

\begin{itemize}
    \item \textbf{Malicious Windows DLL (1st attack):}
        Nel primo attacco, l'attaccante ha inserito un eseguibile malevolo nel tarball,
        e pertanto il numero totale di valori esadecimali nel pacchetto è rimasto invariato, restando a 0.
        Nel grafico \ref{fig:hex_unique_plot2} questi pacchetti sono 
        rappresentati dalle linee colorate tratto e punto.

        Dunque la metrica non è indicativa per rilevare questo specifico attacco.

    \item \textbf{Crypto attack (2nd attack):}
        Nel secondo attacco preso in analisi, l'attaccante ha introdotto una linea di codice 
        offuscata in un file già esistente, aumentando il numero di valori esadecimali unici totali presenti nel pacchetto.
        
        Dal grafico \ref{fig:hex_unique_plot2} si può vedere che 
        queste versioni malevole dei pacchetti (le croci rosse sulle linee colorate tratteggiate),
        hanno un numero di valori esadecimali unici (445) che
        si posizionano tra i valori outlier.
        
        Poichè i valori esadecimali per le versioni malevole si posizionano tra i valori outlier,
        la metrica risulta efficace per rilevare questo specifico attacco.
    
    \item \textbf{Shai-Hulud 1.0 (3rd attack):}   
        Nel terzo attacco, l'attaccante ha introdotto un file malevolo in chiaro, chiamato bundle.js;
        pertanto il numero di valori esadecimali unici non è aumentato in modo significativo.
        
        Dal grafico \ref{fig:hex_unique_plot2} si può vedere che
        il numero di valori esadecimali delle versioni malevole, le croci rosse sulle linee colorate continue, 
        rimangono dentro gli intervalli di confidenza.
        
        Dunque, la metrica non è indicativa per rilevare questo attacco,
        in quanto le versioni malevole sono indistinguibili dalle versioni legittime.
        
    \item \textbf{Shai-Hulud 2.0 (4th attack):}
        Infine, nel quarto attacco preso in analisi, 
        l'attaccante introduce un file malevolo offuscato, chiamato setup\_bun.js.

        Come si può notare nel grafico \ref{fig:hex_unique_plot2},
        i pacchetti con la versione compromessa, ovvero le croci rosse sulle linee colorate puntinate,
        hanno un numero di valori esadecimali unici che schizza fino ad arrivare a 122535; 
        superando di gran lunga i valori massimi outlier.

        Dunque, la metrica risulta efficace per rilevare questo attacco.

\end{itemize}

Ricapitolando, 
la metrica risulta efficente nel rilevare l'attacco crypto (2nd attack) e l'attacco Shai-Hulud 2.0 (4th attack),
in quanto in entrambi i casi, le versioni malevole hanno un numero di valori esadecimali unici tali da
essere valori outlier.
Al contrario, la metrica non è indicativa per rilevare l'attacco Malicious Windows DLL (1st attack) e l'attacco Shai-Hulud 1.0 (3rd attack),
in quanto in entrambi i casi, le versioni malevole hanno un numero di valori esadecimali unici tale da posizionarsi dentro agli intervalli di confidenza,
rendendo le versioni malevole indistinguibili dalle versioni legittime.


\newpage



% ------------------------------------------
\section{Analysis Metric 4: Unique Sensitive keywords (no duplicate values)}

\subsection{Goodware Dataset Scenario}


\subsection{Outlier Analysis}



\subsection{Comparison Malware dataset vs. Goodware dataset and Results for each 4 attacks}

\newpage

% ------------------------------------------
\section{Analysis Metric 5: Unique Cryptocurrency addresses (no duplicate values)}

\subsection{Goodware Dataset Scenario}


\subsection{Outlier Analysis}


\subsection{Comparison Malware dataset vs. Goodware dataset and Results for each 4 attacks}

\newpage























\begin{comment}
Here you will present the results of your work. Your skeptical reader is now asking themself
``is this working?''. You will show here to what extent it does. Be honest about the limitations: if your system doesn't work in some cases, you should say so (and explain where and why, if you
can). Like the \nameref{ch:method} chapter, this one may end up being split into multiple chapters.

Keep in mind that this shouldn't be a mere dump of experimental results; you're rather
\emph{teaching} your reader how and to what extent you managed to solve the problem you
described in the \fullref{ch:introduction}. If you have additional results that may be useful but are not necessary to understand the points you're making (e.g., you evaluated your system on multiple datasets and the results all tell the same story), the place for them is in an appendix.

This chapter should have a lot of links to the methodology chapter(s), because you're
evaluating the choices you made there. If you developed a system made of multiple parts,
make sure that you test them separately and together, so that the reader can understand how
important each part is.

\begin{table}
    \centering
    \begin{tabular}{lrr}
        \toprule
        \textbf{Column 1} & \textbf{Column 2} & \textbf{Column 3} \\
        \midrule
        aaa &   2 &   3.42 \\
        bbb & 505 &   6.00 \\
        ccc &   8 & 901.02 \\
        \bottomrule
    \end{tabular}
    \caption{A table using the \latex \texttt{booktabs} package. Note there are no vertical
    rules. In case you want to do more fancy stuff (e.g., merging cells), check also the
    \texttt{multirow} and \texttt{multicol} packages. It's generally a good idea for readability
    to align text to the left and numbers to the right (use the same number of significant digits).}
    \label{tab:table}
\end{table}

This chapter is likely to be quite full of figures and tables. Try to make them as informative
as possible (e.g., use multiple lines in the same plot if possible). Showing graphs
effectively is a complex art; try to spend some time on it and ask for guidance from your
advisor. Whenever you have a figure or a table, make sure that you refer to it in the text
(after all, if it's not referred to in the text it means it has no part in the story you're
telling, so it has no place in this Chapter). Always use the text (and the caption, if you can) to explain what you want the reader to understand from the figure.

If you are using \LaTeX, it will automatically place figures, tables, algorithms, etc. somewhere in the text.
These parts of the documents are called \emph{floats} because their position floats around depending on \LaTeX's will.
There are options to advise \LaTeX\ about where floats should be placed. My advice is not to waste too much time on this
and trust that they will end up in a good place. Generally, write the floats' code before you refer to them: they will
never be placed before the place they appear in.

For Figures, generally use vectorial formats (e.g., PDF, SVG) where it makes sense if you can.
They will result in higher-quality images that are in most cases also smaller, leading to
shorter compiling times and a smaller resulting PDF.

Tables look better without vertical rules: an example is in \Cref{tab:table}.

\section{Discussion}

At the end of this chapter, take a step back and summarize all the results so that the reader
can understand the big picture. Sometimes, this becomes large and important enough to be worth
a chapter on its own.
\end{comment}